\chapter{Diseño del sistema: agente de coordinación sobre middleware SDK}
\label{chap:implementation-design}

\section{Instanciación del marco}

This chapter instantiates the theoretical framework from Chapter~\ref{chap:theoretical-solution} as a Python coordination agent built over a middleware SDK. The implementation design keeps the theoretical contribution protocol-independent, but it chooses the Agent2Agent (A2A) protocol as the canonical interoperability backbone for concrete communication. A2A is an open protocol intended to support communication and coordination among agents built with different technologies. Its project documentation describes standardized communication, Agent Cards for discovery, flexible interaction patterns, and rich data exchange through text, files, and structured JSON \cite{a2aproject2025}. Google's announcement similarly frames A2A as an open protocol for secure information exchange and coordination among agents across enterprise platforms \cite{google2025a2a}.

The implementation separates two roles. \texttt{CoordinationSdk} is the communication and runtime middleware. It registers remote A2A agents, local Python agents, and Lingo-backed linguistic agents; normalizes their capabilities; dispatches authorized tasks; converts artifacts; records runtime failures; and exposes communication traces. It does not plan, decide feasibility, authorize decompositions, or coordinate complete user requests.

\texttt{CoordinationAgent} is the central coordinating role. It receives the user request, uses the SDK to inspect the registry and capability index, invokes SDK-registered linguistic capabilities when interpretation or candidate planning is needed, builds direct or decomposed candidate plans, calls the symbolic \texttt{FeasibilityAnalyzer}, authorizes or refuses the plan, dispatches only authorized tasks through the SDK, and aggregates the final result. The agent never calls raw A2A transports, raw local handlers, or raw Lingo runtimes directly. All execution traffic still passes through \texttt{CoordinationSdk}.

A2A remains the default wire model for independently deployed agents. Remote agents are admitted through A2A Agent Cards and invoked through A2A-compatible task exchange. Local Python agents and Lingo-based linguistic agents are registered through the SDK and normalized into the same internal agent records, with identity, capabilities, input and output modes, trust metadata, invocation endpoint, and validation contract. Local and linguistic agent authors are therefore not required to implement A2A plumbing, while the coordination agent can reason over a uniform registry.

The prototype also uses Lingo as an optional linguistic and context-engineering substrate \cite{lingo2026}. Lingo is not hidden inside the SDK as a coordinator. Instead, Lingo-backed interpretation, extraction, classification, or candidate-planning behavior is exposed as SDK-registered linguistic capabilities. Non-linguistic local agents can use the SDK without depending on Lingo. The prototype dependency set pins \texttt{lingo-ai==2.0.2} and uses \texttt{a2a-sdk}, \texttt{pydantic}, \texttt{httpx}, \texttt{pytest}, and \texttt{pytest-asyncio} for the executable implementation and tests.

Lingo is therefore an implementation witness for the abstract transitions defined in Chapter~\ref{chap:theoretical-solution}. A Lingo \texttt{Context} can serve as the mutable linguistic ledger used by interpretation and trace generation. \texttt{Engine.create}, \texttt{Engine.decide}, and \texttt{Engine.choose} can implement controlled proposal operations. \texttt{Flow} can encode procedural candidate-planning workflows, while \texttt{StateMachine} can represent deterministic workflow states. Pydantic-backed schemas can materialize capability and artifact contracts, and \texttt{@app.when} reflexes can provide early guardrails or context injection. These mechanisms support interpretation, proposal, and feedback grounding; they do not authorize execution by themselves.

The mapping between formal objects and implementation records is direct. The records defined in Section~\ref{sec:data-structures} materialize the formal objects of Chapter~\ref{chap:theoretical-solution} as follows.

\begin{itemize}
  \item \texttt{CoordinationAgent} materializes the central coordinating authority. It owns interpretation workflow, direct matching, decomposition, feasibility checks, auxiliary decisions, authorization, and final result construction.
  \item \texttt{CoordinationSdk} materializes the communication abstraction. It owns registry admission, A2A gateway behavior, local and Lingo runtimes, task dispatch, artifact normalization, runtime failure recording, and communication trace access.
  \item \texttt{AgentRegistryEntry} materializes an agent $a$ (Definition~\ref{def:agent}), supplying $\mathit{id}$, declared capabilities $C(a)$, input modes $M_{\mathrm{in}}(a)$, output modes $M_{\mathrm{out}}(a)$, local admission policy $\tau(a)$, observable status $\sigma(a)$, and the invocation endpoint used by the SDK.
  \item \texttt{CapabilityRequirement} materializes a required capability descriptor $q \in \mathcal{Q}$ with modalities, schemas, constraints, trust level, and auxiliary eligibility.
  \item \texttt{ProblemRequest} materializes the request $P = (g, \mathcal{Q}, \chi, \mathcal{O}, \Gamma)$ from Definition~\ref{def:problem-request}.
  \item \texttt{SolutionProposal} materializes the coordination plan $\Pi$ from Definition~\ref{def:coordination-plan}, including task requirements, assignments, dependencies, auxiliary specifications, expected artifacts, and completion criteria.
  \item \texttt{FeasibilityReport} records the verdict and predicate evidence from Definition~\ref{def:feasibility}. It is produced by the \texttt{FeasibilityAnalyzer} at the \texttt{CoordinationAgent} layer.
  \item \texttt{GeneratedNlpAgentSpec} materializes the auxiliary capability specification from Definition~\ref{def:auxiliary}, bounded by Rules~\ref{rule:approved-construction}--\ref{rule:lifecycle}.
  \item \texttt{CoordinatorState} is the coordination agent's session state containing the registry snapshot, interpreted request, candidate plan, feasibility evidence, and trace.
  \item \texttt{TraceEvent} records auditable observations across interpretation, authorization, delegation, and artifact handling.
\end{itemize}

Each of the five properties stated in Chapter~\ref{chap:theoretical-solution} (justified acceptance, simplicity, bounded auxiliary synthesis, traceability, and explicit refusal) corresponds to a concrete decision or trace event owned by the \texttt{CoordinationAgent}, supported by communication evidence supplied through the SDK.

\section{Componentes del sistema}

The proposed implementation is organized around a coordinating agent above a middleware SDK. The \texttt{CoordinationAgent} is the planner and authorizer. It receives user requests, maintains coordination state, requests linguistic interpretation when needed, constructs candidate plans, applies deterministic feasibility checks, decides whether to authorize or refuse the plan, dispatches authorized work through the SDK, and aggregates the final result.

\texttt{CoordinationSdk} is the communication and runtime substrate. It is the facade through which agents are registered and invoked. The SDK may use private helper classes internally, but adapter inheritance and transport details are not part of the public API. The SDK can reject invalid configuration, malformed registrations, missing handlers, invalid Agent Cards, and transport-level failures, but it does not decide whether a user request is feasible as a multi-agent plan.

The private SDK subsystems keep communication responsibilities simple. The Registry and Capability Index maintain admitted agent records and searchable capability metadata. The A2A Gateway fetches and validates Agent Cards, sends A2A tasks, and normalizes remote status, artifacts, cancellations, timeouts, and failures. The Local Agent Runtime wraps Python functions, scripts, classes, symbolic planners, and deterministic tools. The Lingo Agent Runtime wraps Lingo apps, flows, tools, state machines, and guardrails when linguistic behavior is required. The SDK Trace Store records communication and runtime observations. The Artifact Normalizer converts outputs into validator-friendly records.

\begin{table}[htbp]
\centering
\caption{Componentes de implementación mapeados a conceptos teóricos.}
\label{tab:implementation-components}
\begin{tabular}{p{0.30\textwidth}p{0.58\textwidth}}
\toprule
\textbf{Componente} & \textbf{Rol en la implementación} \\
\midrule
\texttt{CoordinationAgent} & Posee el flujo de interpretación, correspondencia directa, descomposición, verificaciones de viabilidad, decisiones auxiliares, autorización y construcción del resultado final. \\
\texttt{CoordinationSdk} & Expone operaciones de middleware para registro, inspección del registro, indexación de capacidades, despacho de tareas, normalización de artefactos, reporte de fallos en tiempo de ejecución, trazas y reinicio de sesión. \\
\texttt{FeasibilityAnalyzer} & Aplica verificaciones deterministas para \texttt{CoordinationAgent}; no forma parte de la frontera de comunicación del SDK. \\
Pasarela A2A & Gestiona la ingesta de Agent Cards y la comunicación compatible con A2A con agentes remotos. \\
Runtime de agente local & Envuelve manejadores de Python y herramientas locales deterministas como capacidades de agentes administradas por el SDK. \\
Runtime de agente Lingo & Envuelve aplicaciones, flujos, herramientas, máquinas de estado y guardrails opcionales de Lingo como capacidades ordinarias del SDK. \\
Almacén de trazas del SDK & Rastrea eventos de comunicación, fallos en tiempo de ejecución, estado de tareas y observaciones de artefactos. \\
Normalizador de artefactos & Convierte salidas remotas, locales y lingüísticas en registros utilizables por validadores y agregadores. \\
Fábrica de capacidades auxiliares & Usada por \texttt{CoordinationAgent} para especificar agentes lingüísticos temporales y ligeros cuando están justificados. \\
\bottomrule
\end{tabular}
\end{table}

\section{Registro de agentes y frontera de comunicación}

The SDK admits three kinds of agents through one communication boundary. A remote A2A agent is registered with \texttt{register\_a2a\_agent(agent\_card\_url, trust\_level=...)}. The SDK fetches the Agent Card, validates the metadata, extracts identity, service endpoint, skills, modalities, and security requirements, and converts them into an internal registry entry. Because the external representation remains A2A-compatible, other conforming A2A implementations can participate without custom glue.

A local agent is registered with \texttt{register\_local\_agent(name, capabilities, handler, ...)}. The handler may be a Python function, script wrapper, class, symbolic planner, or deterministic tool. The local agent does not need to know about A2A. The SDK supplies the same identity, capability, modality, authority, invocation, and validation metadata that a remote Agent Card would otherwise provide. During dispatch, the SDK invokes the local handler through the Local Agent Runtime and records the result as a normal task observation.

A linguistic agent is registered through the SDK's linguistic-agent registration method. The wrapped object may be a Lingo app, flow, structured generation routine, state machine, tool set, or guardrail-bearing conversational process. The SDK exposes the linguistic agent as an ordinary capability provider. Lingo remains optional and internal to that capability. The \texttt{CoordinationAgent} may invoke the linguistic capability through \texttt{sdk.send\_task(...)} when it needs interpretation, extraction, classification, explanation, or candidate planning.

No registered agent communicates directly with another registered agent through application-level shortcuts. The \texttt{CoordinationAgent} also does not bypass the SDK to reach remote transports, local handlers, or Lingo runtimes. This preserves one communication boundary while keeping planning and feasibility outside the SDK.

\section{Estructuras de datos mínimas}
\label{sec:data-structures}

The implementation can be described through two public surfaces and a set of typed conceptual records. The records are represented as Pydantic models in the prototype so that Lingo structured generation, deterministic validation, SDK registration, coordination-agent planning, and test fixtures use the same schema surface. The records are not tied to a deployment platform; they define the information that the prototype must preserve.

\begin{Verbatim}[fontsize=\small,breaklines=true,breakanywhere=true]
CoordinationSdk:
  register_a2a_agent(agent_card_url, trust_level="standard")
  register_local_agent(name, capabilities, handler, ...)
  register_linguistic_agent(name, capabilities, lingo_app_or_flow, ...)
  registry_snapshot()
  capability_index()
  send_task(agent_id, task)
  trace()
  reset_session()

CoordinationAgent:
  coordinate(user_request, context=None)
  check_feasibility(problem_or_plan)
  trace()

AgentRegistryEntry:
  agent_id
  name
  agent_kind
  service_endpoint
  invocation_endpoint
  description
  skills
  input_modes
  output_modes
  status
  trust_level
  validation_contract
  source_card

CapabilityRequirement:
  name
  description
  input_schema
  output_schema
  input_modes
  output_modes
  constraints
  auxiliary_eligible
  required_trust_level

ProblemRequest:
  user_goal
  requirements
  constraints
  required_artifacts
  context

TaskSpec:
  task_id
  requirement_name
  assigned_to
  auxiliary_spec_id
  depends_on
  expected_artifacts

FeasibilityReport:
  feasible
  matched_agents
  generated_nlp_agents
  missing_capabilities
  risks
  explanation
  evidence

SolutionProposal:
  tasks
  generated_nlp_agents
  selected_agents
  execution_order
  expected_artifacts
  completion_criteria

GeneratedNlpAgentSpec:
  purpose
  input_schema
  output_schema
  method
  validation_rule
  lifecycle
  authority_bounds
  persists

CoordinatorState:
  registry_snapshot
  interpreted_request
  candidate_plan
  feasibility_evidence
  trace

TraceEvent:
  event_type
  message
  data
  timestamp
\end{Verbatim}

The \texttt{CoordinationSdk} public surface is intentionally limited to registration, inspection, dispatch, trace access, and session control. It hides the differences among remote A2A agents, local Python agents, and Lingo-backed linguistic agents while preserving enough metadata for the coordination agent to apply feasibility rules uniformly. The \texttt{CoordinationAgent} public surface contains the complete coordination workflow and feasibility query. The \texttt{AgentRegistry\allowbreak Entry} materializes the discoverable agent set. The \texttt{Capability\allowbreak Requirement} and \texttt{Problem\allowbreak Request} represent the user's goal, constraints, required capabilities, and required artifacts. The \texttt{Solution\allowbreak Proposal} is a candidate plan, not an executable plan until accepted. The \texttt{Feasibility\allowbreak Report} explains whether the coordination agent can solve the problem and why. The \texttt{Generated\allowbreak Nlp\allowbreak Agent\allowbreak Spec} represents a temporary linguistic capability, such as a classifier or extractor, that can be created for the current solution without becoming a permanent system agent. The \texttt{Coordinator\allowbreak State} and \texttt{Trace\allowbreak Event} records provide the auditable ledger required by the traceability property.

\section{Flujo de trabajo del coordinador}

The application-level coordinator is \texttt{CoordinationAgent}. It talks to the distributed system only through \texttt{CoordinationSdk}, but it keeps planning and authorization outside the SDK. This means the SDK is reusable for non-coordinating clients, while the thesis contribution remains concentrated in the coordination agent's feasibility-aware workflow.

\begin{Verbatim}[fontsize=\small,breaklines=true,breakanywhere=true]
sdk = CoordinationSdk(...)
analyzer = FeasibilityAnalyzer(...)
agent = CoordinationAgent(
    sdk=sdk,
    feasibility_analyzer=analyzer,
    linguistic_interpreter="optional-sdk-agent-id",
    auxiliary_factory=optional_factory,
    artifact_aggregator=aggregator,
)

sdk.register_a2a_agent(agent_card_url, trust_level="standard")
sdk.register_local_agent(name, capabilities, handler, ...)
sdk.register_linguistic_agent(name, capabilities, lingo_app_or_flow, ...)

result = agent.coordinate(user_request, context=session_context)
\end{Verbatim}

Inside \texttt{coordinate}, the \texttt{CoordinationAgent} performs the simple-first decision process. It asks the SDK for the current registry and capability index, transforms or admits the user's message into a typed \texttt{ProblemRequest}, attempts direct capability matching, requests candidate decomposition only when direct delegation is insufficient, checks the proposal with the \texttt{FeasibilityAnalyzer}, adds a bounded auxiliary specification only when the feasibility model permits it, and dispatches authorized tasks through the SDK.

\begin{Verbatim}[fontsize=\small,breaklines=true,breakanywhere=true]
CoordinationAgent.coordinate(problem):
  registry = sdk.registry_snapshot()
  capabilities = sdk.capability_index()
  request = interpret_or_admit_problem(problem, sdk)

  candidate_plan = direct_or_decomposed_plan(request, capabilities, sdk)
  report = feasibility_analyzer.check(request, registry, candidate_plan)

  if report.has_auxiliary_eligible_gap:
      auxiliary = auxiliary_factory.specify(report.gap)
      candidate_plan = candidate_plan.with_auxiliary_agent(auxiliary)
      report = feasibility_analyzer.check(request, registry, candidate_plan)

  if report.is_not_feasible:
      return structured_infeasibility(report)

  results = [
      sdk.send_task(task.assigned_to, task)
      for task in candidate_plan.execution_order
  ]
  return artifact_aggregator.combine(results)
\end{Verbatim}

Task dispatch remains polymorphic inside the SDK but not visible to the coordination agent. If the selected executor is a remote A2A agent, the A2A Gateway sends the task and normalizes the remote response. If it is a local agent, the Local Agent Runtime invokes the registered handler and converts the result into the same task observation shape. If it is a Lingo-backed agent, the Lingo Agent Runtime executes the approved app, flow, tool, or state machine and returns a normalized artifact. In all cases, planning infeasibility is structured data returned by the \texttt{CoordinationAgent}. SDK exceptions are reserved for invalid SDK configuration or genuine transport/runtime failures.

This workflow keeps the LLM useful but bounded. Lingo may supply typed context handling and structured proposal mechanisms through SDK-registered linguistic capabilities, but it does not alone decide feasibility. Deterministic checks remain responsible for explicit properties such as availability, modality compatibility, declared skills, authority, and missing capabilities. Lingo's native tool-calling interface may be useful for controlled experiments that intentionally remove symbolic checks, but the main hybrid implementation uses developer-controlled structured output so that symbolic authorization remains explicit.

\section{Especializaciones de agentes locales y lingüísticos}

The SDK treats local and linguistic agents as specializations of the same communication boundary rather than as separate architectural paths. A local agent is appropriate when the capability is deterministic, symbolic, scriptable, or already implemented in Python. Examples include a rule-based planner, a schema validator, a database wrapper, a file transformer, or a mathematical tool. The SDK supplies the communication wrapper, so the local agent author focuses on the capability contract and handler behavior.

A linguistic agent is appropriate when the capability transforms, interprets, classifies, summarizes, or produces language-mediated artifacts. It may be implemented with Lingo, an LLM, a smaller NLP model, rules, retrieval, or a hybrid pipeline. When a registered linguistic capability is implemented with Lingo, it follows the primitive mapping described at the beginning of this chapter; other implementations must still expose the same typed capability, artifact, and validation contracts to the SDK.

Generated NLP agents are solution-local. The \texttt{CoordinationAgent} includes their specifications in the \texttt{SolutionProposal}, including purpose, method, expected input, expected output, validation rule, and lifecycle. They do not automatically enter the registry. Promotion to a persistent registered agent would require a separate SDK admission process.

\section{Diagrama de arquitectura}

Figure~\ref{fig:implementation-architecture} summarizes the implementation structure. The diagram is intentionally abstract enough to support later coding choices while still showing that coordination logic belongs to \texttt{CoordinationAgent}, while communication and runtime adaptation belong to \texttt{CoordinationSdk}.

\begin{figure}[htbp]
\centering
\fbox{
\begin{minipage}{0.90\textwidth}
\centering
Problema del usuario o solicitud de aplicación\\
$\downarrow$\\
\texttt{CoordinationAgent}\\
$\downarrow$\\
\begin{tabular}{ccc}
Interpretación & Viabilidad & Agregación \\
\end{tabular}\\
$\downarrow$\\
Frontera de middleware \texttt{CoordinationSdk}\\
$\downarrow$\\
\begin{tabular}{ccc}
Pasarela A2A & Runtime local & Runtime Lingo \\
\end{tabular}\\
$\downarrow$\\
\begin{tabular}{ccc}
A2A remotos & Python locales & Lingüísticos \\
\end{tabular}\\
$\downarrow$\\
Artefactos, estados, fallos y traza de comunicación\\
$\downarrow$\\
Resultado final o informe estructurado de inviabilidad
\end{minipage}
}
\caption{\texttt{CoordinationAgent} sobre middleware SDK para coordinación compatible con A2A entre agentes remotos, locales y lingüísticos.}
\label{fig:implementation-architecture}
\end{figure}

\section{Secuencia de ejecución}

The execution sequence begins when the user or application submits a problem request to \texttt{CoordinationAgent.coordinate}. The coordination agent asks the SDK for the current registry snapshot and capability index. It obtains a typed problem representation from a structured caller or by invoking a registered linguistic capability through the SDK. The coordination agent then attempts direct matching before considering decomposition. The \texttt{FeasibilityAnalyzer} checks the candidate plan and either approves it, permits an auxiliary NLP specification to be added, or rejects it as infeasible. Approved tasks are dispatched through \texttt{sdk.send\_task(...)}. The SDK chooses the correct internal runtime, records status and failures, normalizes artifacts, and exposes communication trace events. The coordination agent aggregates the authorized outputs or returns an infeasibility report.

\begin{figure}[htbp]
\centering
\fbox{
\begin{minipage}{0.88\textwidth}
\begin{enumerate}
\item El usuario o la aplicación envía un problema y restricciones a \texttt{CoordinationAgent}.
\item El agente de coordinación lee el registro y las capacidades mediante \texttt{CoordinationSdk}.
\item El agente de coordinación obtiene o construye un \texttt{ProblemRequest} tipado.
\item El agente de coordinación intenta correspondencia directa de capacidades antes de la descomposición.
\item El agente de coordinación solicita planificación candidata solo cuando la correspondencia directa es insuficiente.
\item El analizador de viabilidad valida habilidades, modalidades, dependencias, autoridad y contratos de evidencia.
\item La fábrica de capacidades auxiliares añade una especificación temporal solo cuando está justificada.
\item El agente de coordinación despacha tareas autorizadas mediante \texttt{sdk.send\_task(...)}.
\item El SDK registra estado en tiempo de ejecución, errores, artefactos y evidencia de comunicación.
\item El agente de coordinación devuelve el resultado final o un informe estructurado de inviabilidad.
\end{enumerate}
\end{minipage}
}
\caption{Secuencia de ejecución para la resolución de problemas mediada por \texttt{CoordinationAgent}.}
\label{fig:execution-sequence}
\end{figure}

\section{Manejo de fallos e inviabilidad}

The implementation treats infeasibility as a first-class coordination-agent result. If no remote, local, linguistic, or approved auxiliary capability can satisfy a required subtask, the \texttt{CoordinationAgent} returns a structured report. If an agent is unavailable, the coordination agent may select another compatible agent from the SDK registry; if no replacement exists, it reports the missing dependency. If modalities are incompatible, the report identifies which input or output mode cannot be bridged. If a linguistic planner proposes a decomposition that violates explicit constraints, the \texttt{FeasibilityAnalyzer} rejects it.

The SDK normalizes runtime failures rather than planning failures. A2A transport failures, local handler exceptions, Lingo execution failures, invalid artifacts, cancellations, and timeouts are recorded as structured observations linked to the relevant task and agent. This behavior is essential for simplicity and trust. The system should not create complex plans merely to appear capable. It should coordinate only when the available and generated capabilities justify the solution. Otherwise, the coordination agent should explain the boundary clearly. Lingo guardrails may detect obvious unsafe, irrelevant, or ambiguous inputs early, but such reflexes are treated as advisory preprocessing rather than as replacements for the feasibility verdict.

\section{Síntesis del capítulo}

The design specifies a distributed multi-agent environment with a clear split between coordination and communication. \texttt{CoordinationAgent} owns planning, feasibility, authorization, auxiliary decisions, and aggregation. \texttt{CoordinationSdk} owns A2A interoperability, local runtime invocation, Lingo runtime invocation, registry normalization, task dispatch, artifact normalization, communication failures, and trace access. Agents remain independently deployable and vendor independent. A2A provides the interoperability backbone through Agent Cards and task exchange. Local Python agents and Lingo-backed linguistic agents are wrapped by the SDK so that they participate through the same registry, dispatch, artifact, and trace mechanisms as remote A2A agents. The evaluation defined in Chapters~\ref{chap:methodology} and~\ref{chap:results} must determine whether an implementation preserves both communication-boundary compliance and coordination-ownership compliance in practice.
